\documentclass[12pt]{article}

\usepackage{amsmath}
\usepackage{geometry}
\geometry{margin=1in}

\setlength{\parskip}{0.5em}

\begin{document}

\title{Kinematics: Projection Practice 1}
\author{}
\date{}
\maketitle

\vspace{1em}

\textbf{Instructions:}

\begin{enumerate}
\item Describe the motion before writing formulas.
\item Show all steps clearly.
\item Include units in every step and in the final answer.
\end{enumerate}

\noindent\rule{\textwidth}{0.4pt}

\newpage

\section*{Projectile Motion Practice Problems}

\subsection*{1. Components of Horizontal Projectile Motion}
An object undergoing horizontal projectile motion performs \underline{\hspace{3cm}} motion in the horizontal direction and \underline{\hspace{3cm}} motion in the vertical direction. If the initial horizontal velocity is increased, the time the object remains in the air will \underline{\hspace{3cm}} (choose from: "increase", "decrease", or "remain unchanged").

\subsection*{2. Horizontal Launch from a Platform}
A small ball is thrown horizontally from the edge of a platform at a height of $19.6 \text{ m}$ with an initial velocity of $10 \text{ m/s}$. Neglecting air resistance, take the acceleration due to gravity $g = 9.8 \text{ m/s}^2$.

\textbf{Find:} The time taken for the ball to hit the ground and the horizontal distance from the base of the platform to the landing point.
\newpage
\subsection*{3. Targeting a Package Drop}
An airplane is flying at a constant horizontal speed of $100 \text{ m/s}$ at an altitude of $2000 \text{ m}$. In order for a dropped package to hit a specific target on the ground, at what horizontal distance from the target should the pilot release the package? (Take $g = 10 \text{ m/s}^2$)

\newpage
\subsection*{4. Impact Angle and Initial Height}
An object performs horizontal projectile motion. Upon hitting the ground, the direction of its velocity makes a $45^\circ$ angle with the horizontal. Given that the initial velocity is $20 \text{ m/s}$, calculate the height from which the object was thrown. (Take $g = 10 \text{ m/s}^2$)

\vspace{15cm}

\subsection*{5. Velocity and Acceleration at Peak Height}
An object is launched with an initial velocity $v_0 = 30 \text{ m/s}$ at an angle of $30^\circ$ above the horizontal. At the highest point of its trajectory, the magnitude of its instantaneous velocity is \underline{\hspace{2cm}} $\text{ m/s}$, and the magnitude of its acceleration is \underline{\hspace{2cm}} $\text{ m/s}^2$. (Take $g = 10 \text{ m/s}^2$ and $\cos 30^\circ = \frac{\sqrt{3}}{2}$)
\end{document}